%------------------------------------------------------------------------------%

\usepackage{calculo_varias_variaveis-1}

\title {Derivadas Parciais}

%------------------------------------------------------------------------------%
\begin{document}

\addtitle

%------------------------------------------------------------------------------%
\section{Derivadas Parciais}

%------------------------------------------------------------------------------%
\begin{frame}{Derivada Ordinária -- Uma variável}

  Variação da função com relação a variável

  \pause
  \vspace{\baselineskip}
  Derivada da função \(f\colon R\subset\R\to\R\) no ponto $x_0$

  \[
    \frac{df}{dx}\evalat{x=x_0}{}
    =\; \lim_{h\to 0} \frac{f(x_0+h)-f(x_0)}{h}
  \]
\end{frame}

%------------------------------------------------------------------------------%
\begin{frame}{Derivada Parcial}
  Variação da função com relação a \emph{uma} das variáveis

  considerando as demais \emph{fixas}
\end{frame}

%------------------------------------------------------------------------------%
\begin{frame}{Definição -- Derivada Parcial em $x$}
  A derivada parcial, com relação a \emph{$x$},
  da função \(f\colon R\subset\R^2\to\R\), \(f(x, y)\)

  no ponto \((x, y) = (x_0, y_0)\)

  \[
    \frac{\partial f}{\partial x}\evalat{(x, y) = (x_0, y_0)}{}
    =\; \lim_{h\to 0} \frac{f(\emph{x_0+h}, y_0)-f(x_0, y_0)}{h}
  \]

  \vspace{\baselineskip}
  Desde que o limite exista
\end{frame}

%------------------------------------------------------------------------------%
\begin{frame}{Definição -- Derivada Parcial em $y$}
  A derivada parcial, com relação a \emph{$y$},
  da função \(f\colon R\subset\R^2\to\R\), \(f(x, y)\)

  no ponto \((x, y) = (x_0, y_0)\)

  \[
    \frac{\partial f}{\partial y}\evalat{(x, y) = (x_0, y_0)}{}
    =\; \lim_{h\to 0} \frac{f(x_0, \emph{y_0+h})-f(x_0, y_0)}{h}
  \]

  \vspace{\baselineskip}
  Desde que o limite exista
\end{frame}

%------------------------------------------------------------------------------%
\begin{frame}{Notações}
  \[
    \frac{\partial f}{\partial y}(x, y) \qquad\qquad
    f_y(x, y)
  \]

  \vspace{\baselineskip}
  \[
    \frac{\partial f}{\partial y} \qquad\qquad\qquad
    f_y
  \]
\end{frame}

%------------------------------------------------------------------------------%
\framedgraphic{Retas Tangentes}{retas_tangentes}{}

\begin{frame}{Retas Tangentes}
  Temos duas retas tangentes, uma na direção $x$ e uma na direção $y$

  \pause
  \vspace{\baselineskip}
  Nem sempre o plano gerado por elas é plano tangente da superfície
\end{frame}

%------------------------------------------------------------------------------%
\section{Exemplos}

%------------------------------------------------------------------------------%
\addtocounter{exemplo}{1}
\begin{frame}{Exemplo \theexemplo}
  Encontre os valores de \(\D{f}{x}\) e \(\D{f}{y}\) no ponto \((4, -5)\)
  para a função
  \[
    f(x, y) = x^2 + 3xy + y - 1
  \]
\end{frame}

%------------------------------------------------------------------------------%
\begin{frame}{Exemplo \theexemplo}
  \begin{align*}
    \onslide<+->{\D{f}{x}}
    \onslide<+->{& = \D{}{x} \left(x^2 + 3xy + y - 1\right)} \\[2mm]
    \onslide<+->{& = \D{x^2}{x} + \D{3xy}{x} + \D{y}{x} - \D{1}{x}} \\[2mm]
    \onslide<+->{& = 2x + 3y\D{x}{x} + 0 - 0} \\[2mm]
    \onslide<+->{& = 2x + 3y}
  \end{align*}
\end{frame}

%------------------------------------------------------------------------------%
\begin{frame}{Exemplo \theexemplo}
  \[
    \D{f}{x}(4, -5)
    \pause
    = \left(2x + 3y\right)\evalat{(x, y)=(4, -5)}{}
    \pause
    = 2\cdot 4 + 3(-5)
    \pause
    = -7
  \]
\end{frame}

%------------------------------------------------------------------------------%
\begin{frame}{Exemplo \theexemplo}
  \begin{align*}
    \onslide<+->{\D{f}{y}}
    \onslide<+->{& = \D{}{y} \left(x^2 + 3xy + y - 1\right)} \\[2mm]
    \onslide<+->{& = \D{x^2}{y} + \D{3xy}{y} + \D{y}{y} - \D{1}{y}} \\[2mm]
    \onslide<+->{& = 0 + 3x\D{y}{y} + \D{y}{y} - 0} \\[2mm]
    \onslide<+->{& = 3x + 1}
    \onslide<+->{\hspace{60mm}\text{Cadê o $y$?}}
  \end{align*}
\end{frame}

%------------------------------------------------------------------------------%
\begin{frame}{Exemplo \theexemplo}
  \[
    \D{f}{y}(4, -5)
    \pause
    = \left(3x + 1\right)\evalat{(x, y)=(4, -5)}{}
    \pause
    = 3\cdot 4 + 1
    \pause
    = 13
  \]
\end{frame}

%------------------------------------------------------------------------------%
\addtocounter{exemplo}{1}
\begin{frame}{Exemplo \theexemplo}
  Encontre \(\D{f}{y}\) para \(f(x, y) = y\sin(xy)\)
\end{frame}

%%------------------------------------------------------------------------------%
%\begin{frame}{Exemplo \theexemplo}
%  \begin{align*}
%    \D{f}{x}
%    & = \D{}{x} \big(y\sin(xy)\big) \\[2mm]
%    & = y \D{}{x}\sin(xy)  \\[2mm]
%    & = y y \cos(xy)  \\[2mm]
%    & = y^2 \cos(xy)
%  \end{align*}
%\end{frame}

%------------------------------------------------------------------------------%
\begin{frame}{Exemplo \theexemplo}
\begin{align*}
  \onslide<+->{\D{f}{y}}
  \onslide<+->{& = \D{}{y} \big(y\sin(xy)\big)} \\[2mm]
  \onslide<+->{& = \D{y}{y} \sin(xy) + y \D{}{y}\sin(xy)} \\[2mm]
  \onslide<+->{& = \sin(xy) + y x\cos(xy)}
\end{align*}
\end{frame}

%------------------------------------------------------------------------------%
\addtocounter{exemplo}{1}
\begin{frame}{Exemplo \theexemplo}
  Encontre as funções $f_x$ e $f_y$ sendo que
  \(
    f(x, y) = \frac{2y}{y+\cos(x)}
  \)
\end{frame}

%------------------------------------------------------------------------------%
\begin{frame}{Exemplo \theexemplo}
\begin{align*}
  \onslide<+->{\D{f}{x}}
  \onslide<+->{& = 2\;\D{}{x} \left(\frac{y}{y+\cos(x)}\right)} \\[2mm]
  \onslide<+->{& = 2\;\dfrac{\displaystyle\;\D{y}{x}\left(y+\cos(x)\right) - y \D{}{x}\left(y+\cos(x)\right)\;}
               {\left(y+\cos(x)\right)^2}} \\[2mm]
  \onslide<+->{& = 2\;\dfrac{\; 0 - y \left(0-\sin(x)\right)\;}
               {\left(y+\cos(x)\right)^2}} \\[2mm]
  \onslide<+->{& = \frac{2y \sin(x)\;}{\left(y+\cos(x)\right)^2}}
\end{align*}
\end{frame}

%------------------------------------------------------------------------------%
\begin{frame}{Exemplo \theexemplo}
\begin{align*}
  \onslide<+->{\D{f}{y}}
  \onslide<+->{& = 2\;\D{}{y} \left(\frac{y}{y+\cos(x)}\right)} \\[2mm]
  \onslide<+->{& = 2\;\dfrac{\displaystyle\;\D{y}{y}\left(y+\cos(x)\right) - y \D{}{y}\left(y+\cos(x)\right)\;}
                            {\left(y+\cos(x)\right)^2}} \\[2mm]
  \onslide<+->{& = 2\;\dfrac{\;\left(y +\cos(x)\right) - y(1 + 0)\;}
                            {\left(y+\cos(x)\right)^2}} \\[2mm]
  \onslide<+->{& = \frac{2\cos(x)}{\,\left(y+\cos(x)\right)^2\,}}
\end{align*}
\end{frame}

%------------------------------------------------------------------------------%
\addtocounter{exemplo}{1}
\begin{frame}{Exemplo \theexemplo}
  Encontre \(\D{z}{x}\) se a equação
  \[
    yz - \ln(z) = x + 1
  \]
  define $z$ como uma função de $x$ e $y$

  \pause
  \vspace{\baselineskip}
  Vamos derivar os dois lados por $x$, assumindo \(z=z(x, y)\)
\end{frame}

%------------------------------------------------------------------------------%
\begin{frame}{Exemplo \theexemplo}
  \begin{align*}
    \onslide<+->{yz - \ln(z) & = x + 1} \\[2mm]
    \onslide<+->{\D{}{x}\big(yz - \ln(z)\big) & = \D{}{x}\left(x + 1\right)} \\[2mm]
    \onslide<+->{\D{}{x}(yz) - \D{}{x}\big(\ln(z)\big) & = 1} \\[2mm]
    \onslide<+->{y\D{z}{x} - \frac{1}{z}\D{z}{x} & = 1} \\[2mm]
    \onslide<+->{\D{z}{x}\left( y - \frac{1}{z} \right) & = 1}  \\[2mm]
    \onslide<+->{\D{z}{x}
    & = \left(y - \frac{1}{z}\right)^{-1}}
    \onslide<+->{  = \left(\frac{yz-1}{z}\right)^{-1}}
    \onslide<+->{  =  \frac{z}{yz-1}}
  \end{align*}
\end{frame}

%------------------------------------------------------------------------------%
\section{Derivadas Parciais e Continuidade}

%------------------------------------------------------------------------------%
\begin{frame}{Derivada e Continuidade em Uma Variável}
  No cálculo de funções de \emph{uma variável}

  \pause
  se uma função é derivável em um ponto $x_0$

  \pause
  ela é contínua nesse ponto
\end{frame}

%------------------------------------------------------------------------------%
\begin{frame}{Derivadas Parciais e Continuidade}
  Se as derivadas parciais de $f$ existem em $(x_0, y_0)$

  \pause
  \emph{não podemos dizer nada} sobre a continuidade de $f$ no ponto
\end{frame}

%------------------------------------------------------------------------------%
\addtocounter{exemplo}{1}
\begin{frame}{Exemplo \theexemplo}
  Considerando a função
  \(
  f(x, y) = \begin{cases}
    0, & xy \neq 0 \\
    1, & xy = 0
  \end{cases}
  \)
  \begin{enumerate}
    \item Prove que $f$ não é contínua na origem
    \item Mostre que suas derivadas parciais existem na origem
  \end{enumerate}
\end{frame}

%------------------------------------------------------------------------------%
\begin{frame}{Exemplo \theexemplo}
  1 -- Queremos mostrar que $f$ não é contínua na origem

  \pause
  Verificando a condição 2: existência do limite na origem

  \pause
  Considerando a reta $y=0$, com $x\neq 0$
  \pause
  \[
    f(x, y)\evalat{y=0}{}
    \pause
    = f(x, 0)
    \pause
    = 1
  \]

  \pause
  Considerando a reta $y=x$, com $x\neq 0$
  \[
    \pause
    f(x, y)\evalat{y=x}{}
    \pause
    = f(x, x)
    \pause
    = 0
  \]

  \pause
  Como $f$ tende a valores diferentes por
  caminhos diferentes para a origem,\\
  o limite não existe
\end{frame}

\framedgraphic{Derivada e Continuidade}{Derivadas_parciais_continuidade}{}

%------------------------------------------------------------------------------%
\section{Lista Mínima}

%------------------------------------------------------------------------------%
\begin{frame}{Lista Mínima}
  Cálculo Vol. 2 do Thomas 12\textsuperscript{a} ed. -- Seção 14.3
  \begin{enumerate}
    \item Estudar o texto da seção
    \item Resolver os exercícios: 4, 10, 37-40
  \end{enumerate}

  \vfill
  Atenção: A prova é baseada no livro, não nas apresentações
\end{frame}

%------------------------------------------------------------------------------%
\end{document}
%------------------------------------------------------------------------------%
